
\def\PUDcodigo{ACE014}

\if\COMPILAR0
   \def\PUDcodacad{04507.19}
\else
   \def\PUDcodacad{04506.19}
\fi

\def\PUDdisciplina{Metodologia Científica e Tecnológica}

\def\PUDsemestre{S4}

\def\PUDhoras{40}

%NOVOS ---------------------------------------------------
\def\PUDhorasTeoria{28}
\def\PUDhorasPratica{12}

\if\COMPILAR0
   \def\PUDinterECA{ \hyperref[PUD:ACE020]{TCC (04507.48)} }
\else
   \def\PUDinterECA{ \hyperref[PUD:ACE020]{TCC (04506.47)} }
\fi

\def\PUDatuaisECA{---}

\def\PUDrecursos{Os seguintes recursos poderão ser utilizados: quadro e pinceis; projetor de Multimídia; material impresso; computadores do laboratório de informática do curso.}

%----------------------------------------------------------

\def\PUDcreditos{2}

\def\PUDprerequisito{---}

\def\PUDpreacad{---}

\def\PUDobjetivos{Compreender os fundamentos de metodologia científica, bem como a comunicação científica. Analisar  gêneros acadêmico-científicos. Produzir pré-projeto e projeto de pesquisa.
%Conhecer e correlacionar os fundamentos, os métodos e as técnicas de análise presentes na produção do conhecimento científico;  Compreender as diversas fases de elaboração e desenvolvimento de pesquisas e trabalhos acadêmicos;  Elaborar e desenvolver pesquisas e trabalhos científicos.
}

\def\PUDementa{Fundamentos de Metodologia Científica. Comunicação Científica. Comunicação entre orientandos/orientadores. Tipos de Conhecimento e Ciência. Métodos, Técnicas e Procedimentos de pesquisa científica. Etapas da pesquisa científica. Análise da estrutura e elaboração de gêneros acadêmico-científicos, segundo o Manual do IFCE. Produção do projeto de pesquisa referente aos assuntos vistos no curso.
%Métodos de pesquisa para o trabalho de reflexão, enfatizando os nexos entre conhecimento técnico-profissional e outros campos de conhecimento, como as humanidades;  Sistemática geral da pesquisa, focalizando a definição do objetivo/problema, da contextualização teórica e elaboração de uma proposta de trabalho;  Técnicas de coleta, sistematização, análise e apresentação de informações;  Elaboração e desenvolvimento de pesquisas e trabalhos científicos obedecendo às orientações e normas vigentes nas Instituições de Ensino e Pesquisa no Brasil e na Associação Brasileira de Normas Técnicas.
}

\def\PUDconteudo{ & Metodologia Científica: Fundamentos de metodologia científica e comunicação científica. Tipos de conhecimentos e Ciência. Ciência: conceito e classificação.
 \\
& Pesquisa Científica: Pesquisa científica - conceituação e tipos.	Métodos, técnicas e procedimentos da pesquisa. Etapas da pesquisa. 
\\
& Trabalho Acadêmico-Científicos: Tipos de leitura e esquematização. Normas para elaboração de trabalhos acadêmico-científicos. Editoração: linguagem científica, citações, notas de rodapé, referências bibliográficas, aspectos formais. Fichamentos, sínteses, resumos e resenhas. 
\\
& Produção da Pesquisa Acadêmica: A comunicação e o papel de orientando/orientador. Pré-projeto e Projeto de pesquisa: definição e estrutura.
%&  Fundamentos da Metodologia Científica: Definições conceituais.  Valores e ética no processo de pesquisa;  
%\\ 
% &  A comunicação Científica: O sistema de comunicação na ciência: canais informais e canais formais;  
%\\ 
% &  Métodos e técnicas de pesquisa: Tipos de conhecimento.  Tipos de Ciência.  Classificação das Pesquisas Científicas.  A necessidade e os tipos do Método.  As etapas da pesquisa; 
%\\ 
% &  A comunicação entre orientados/orientadores: O papel de orientado/orientador na produção da pesquisa acadêmica; 
%\\ 
% &  Normas para Elaboração de Trabalhos Acadêmicos: Estrutura e Definição; 
%\\ 
% &  O pré-projeto de pesquisa: Definição.  Modelos.  Elementos; 
%\\ 
% &  O projeto de pesquisa: Definição.  Modelos.  Elementos; 
%\\ 
% &  A organização de texto científico (normas ABNT): Normas para elaboração de trabalhos acadêmicos. 
}

\def\PUDmetodo{A aula será expositiva e dialógica com aplicação de exercícios de forma individual e/ou em pequenos grupos. Apresentação de seminário. Leitura, análise e elaboração de trabalhos científicos. Prática individual em laboratório de escrita.}

\def\PUDavalia{Acompanhamento contínuo do discente por meio de instrumentos e técnicas diversificadas de avaliação que tenham objetivos e critérios bem explicitados. Os aspectos quantitativos da avaliação ocorrerão de acordo com o Regulamento da Organização Didática (ROD) do IFCE.}

\def\PUDbibbasica{ & \href{www.biblioteca.ifce.edu.br/index.asp?codigo_sophia=61981}{Almeida}, Mário de Souza.  Elaboração de projeto, TCC, dissertação e tese: uma abordagem simples, prática e objetiva.  2º ed.  São Paulo, SP: Atlas, 2014.  82p.  ISBN: 9788522491155.
\\
 & \href{www.biblioteca.ifce.edu.br/index.asp?codigo_sophia=24139}{Marconi}, Marina de Andrade.  Fundamentos de metodologia científica.  7º ed.  São Paulo, SP: Atlas, 2010.  297p.  ISBN: 9788522457588.
\\
& \href{www.biblioteca.ifce.edu.br/index.asp?codigo_sophia=23788}{Severino}, Antônio Joaquim.  Metodologia do trabalho científico.  23º ed.  São Paulo, SP: Cortez, 2009.  304p.  ISBN: 9788524913112.
}

\def\PUDbibcomplementar{ & \href{biblioteca.ifce.edu.br/index.asp?codigo_sophia=24239}{Andrade}, Maria Margarida de. Introdução à metodologia do trabalho científico: elaboração de trabalhos na graduação. 10. ed. São Paulo: Atlas, 2010. 158 p. Inclui bibliografia e índice. ISBN 9788522458561.
\\
& \href{www.biblioteca.ifce.edu.br/index.asp?codigo_sophia=61856}{Cervo}, Amado Luiz.  Metodologia científica.  6º ed.  São Paulo, SP: Pearson Prentice Hall, 2007.  162p.  ISBN: 9788576050476. 
 \\
 & \href{www.biblioteca.ifce.edu.br/index.asp?codigo_sophia=61795}{Demo}, Pedro.  Metodologia científica em ciências sociais.  3º ed.  São Paulo, SP: Atlas, 1995.  293p.  ISBN: 9788522412419.
\\
& \href{biblioteca.ifce.edu.br/index.asp?codigo_sophia=61868}{Iskandar}, Jamil Ibrahim. Normas da ABNT: comentadas para trabalhos científicos. 6. ed. Curitiba: Juruá, 2016. 98 p. ISBN 9788536258591.
\\
 %& \href{www.biblioteca.ifce.edu.br/index.asp?codigo_sophia=65145}{Madureira}, Omar Moore de.  Metodologia do projeto: planejamento, execução e gerenciamento.  2º ed.  São Paulo, SP: Blucher, 2015.  361p.  ISBN: 9788521209133.
%\\ 
 & \href{www.biblioteca.ifce.edu.br/index.asp?codigo_sophia=62418}{Koche}, José Carlos.  Fundamentos de metodologia científica: teoria da ciência e iniciação à pesquisa.  34º ed.  Rio de Janeiro, RJ: Vozes, 2015.  182p.  ISBN: 9788532618047. 
\\
 & \href{www.biblioteca.ifce.edu.br/index.asp?codigo_sophia=61933}{Salomon}, Délcio Vieira. Como fazer uma monografia.  13º ed.  São Paulo, SP: WMF Martins Fontes, 2014.  425p.  ISBN: 9788578279004.
 }

\def\PUDautor{Fabrício Bandeira da Silva}

\def\PUDdata{2013-04-25}

\def\PUDrevisao{3}

\def\PUDrevisor{Maria do Socorro Cardoso de Abreu}

\def\PUDdatarevisao{2019-05-15}

\PUDcorpo
